\section{Discussion and Limitations}
\label{sec:discussion}

The most general lesson of this project is that matched-budget kernel tuning is less about finding any fast schedule than about deciding which evidence is worth paying for. Compile-adjacent signals are clearly useful, but primarily as a way to reject obviously poor candidates. Once the set of remaining candidates has narrowed to several plausible schedule families, the same signals are too weak to solve the ranking problem on their own. We believe this observation generalizes more easily than any specific speedup number in this paper, because it explains where lightweight tuning should and should not be trusted.

The second lesson is that workload design can change the apparent quality of a selector. The aligned matrix multiplication results are kept in the paper precisely because they show how a convenient workload can make a selector look stronger than it is on representative shapes. The strongest methodological claim of the study is therefore not that one selector outperforms another on a single benchmark, but that representative workloads are a necessary precondition for any such comparison to be scientifically meaningful.

The third lesson is that bounded negative results improve the final story. Layer normalization is deliberately kept as a secondary, regime-split result because the evidence never justified a strong matched-budget success claim for it. The reduction-splitting parameter for matrix multiplication was removed from the final evaluation surface because it collapsed on the enlarged search space. These outcomes are not detours around the main contribution; they are the reason the final non-reduction-split result can be presented credibly rather than as the survivor of uncontrolled heuristic growth.

The scope of this paper is intentionally narrower than what the broader autotuning literature often targets. We show that a fixed-kernel Triton tuning study, built around an explicit evidence hierarchy and an explicit measurement budget, can produce interpretable lessons about pruning, ranking, workload realism, profiling, and transfer. A universal Triton autotuner, a multi-architecture result, or a substitute for richer compiler-level schedule synthesis are all outside this scope.

\subsection{Threats to validity}

\begin{itemize}
\item \textbf{Single architecture.} All primary studies use a single homogeneous pool of NVIDIA RTX A6000 GPUs. We make no claims about how the selector would behave on newer or substantially different GPU architectures.
\item \textbf{Fixed kernel implementations.} The Triton implementation of each kernel is held constant throughout the study. Broader algorithmic alternatives, kernel fusion, and full-kernel synthesis are outside our scope.
\item \textbf{Bounded schedule spaces.} Every reported comparison comes from a schedule space with explicit admissibility bounds. This improves fairness and interpretability but limits the generality of any single speedup number.
\item \textbf{Limited kernel coverage.} Matrix multiplication is the primary workload and layer normalization is secondary. Whether the same lessons extend to all Triton kernels, particularly those with very different compute-to-memory profiles or irregular memory access patterns, remains an open question.
\item \textbf{No head-to-head comparison with existing autotuners.} We compare selectors within our own framework under a matched budget, but we do not benchmark against external autotuning systems such as Kernel Tuner or AutoTVM. Such a comparison would require matching not only the measurement budget but also the implementation and search strategy, which is a separate study.
\item \textbf{Small effect size.} The final headline speedup of $1.03\times$ is small. While it is stable on repeated runs at a fixed seed, a cross-seed robustness sweep with more seeds would strengthen the claim. We report the improvement as bounded for this reason.
\item \textbf{Empirical rather than theoretical scope.} The contribution is a careful empirical study, not a formal characterization of the schedule selection problem or of its optimal budget allocation.
\end{itemize}

These limitations are part of the paper's argument rather than an afterthought: the value of the study lies precisely in stating where the method works, where it falls short, and why the final positive claim is intentionally bounded.
